% =====================================================================
%  INT210 - Machine Learning aplicado a la Ciberseguridad y al
%  Diseño/Animación Digital
%  Semana 1: Fundamentos de IA, ML y la Jerarquía de Necesidades de Datos
%  Compilar con:  tectonic semana-1-slides.tex   (o xelatex / lualatex)
% =====================================================================
\documentclass[10pt,aspectratio=169]{beamer}

% ---------------------------------------------------------------- Idioma
\usepackage{fontspec}
\usepackage[spanish,es-noshorthands,es-tabla]{babel}

% ---------------------------------------------------------------- Paquetes
\usepackage{amsmath,amssymb,amsfonts}
\usepackage{booktabs}
\usepackage{array}
\usepackage{tikz}
\usetikzlibrary{shapes.geometric,positioning,arrows.meta,calc,fit}
\usepackage{pgfplots}
\pgfplotsset{compat=1.18}
\usepackage{listings}
\usepackage{xcolor}

% ---------------------------------------------------------------- Tema
\usetheme{Madrid}
\usecolortheme{default}
\definecolor{usfxazul}{RGB}{0,51,102}
\definecolor{usfxrojo}{RGB}{153,0,0}
\definecolor{usfxdorado}{RGB}{196,145,2}
\definecolor{codigofondo}{RGB}{246,248,250}
\definecolor{codigoverde}{RGB}{0,128,0}
\definecolor{codigogris}{RGB}{110,110,110}
\definecolor{vp}{RGB}{46,125,50}
\definecolor{fp}{RGB}{239,108,0}
\definecolor{fn}{RGB}{198,40,40}

\setbeamercolor{structure}{fg=usfxazul}
\setbeamercolor{palette primary}{bg=usfxazul,fg=white}
\setbeamercolor{palette secondary}{bg=usfxazul!85!black,fg=white}
\setbeamercolor{palette tertiary}{bg=usfxazul!70!black,fg=white}
\setbeamercolor{frametitle}{bg=usfxazul,fg=white}
\setbeamercolor{title}{bg=usfxazul,fg=white}
\setbeamercolor{block title}{bg=usfxazul,fg=white}
\setbeamercolor{block body}{bg=usfxazul!6}
\setbeamercolor{block title example}{bg=vp,fg=white}
\setbeamercolor{block body example}{bg=vp!6}
\setbeamercolor{block title alerted}{bg=usfxrojo,fg=white}
\setbeamercolor{block body alerted}{bg=usfxrojo!6}
\setbeamertemplate{navigation symbols}{}
\setbeamertemplate{itemize items}[circle]
\setbeamertemplate{enumerate items}[default]
\setbeamertemplate{blocks}[rounded][shadow=false]

% Diapositiva de sección automática
\AtBeginSection[]{
  \begin{frame}{Contenido}
    \tableofcontents[currentsection,hideothersubsections]
  \end{frame}
}

% ---------------------------------------------------------------- Código
\lstdefinestyle{python}{
  language=Python,
  basicstyle=\ttfamily\scriptsize,
  keywordstyle=\color{usfxazul}\bfseries,
  stringstyle=\color{usfxrojo},
  commentstyle=\color{codigoverde}\itshape,
  numberstyle=\tiny\color{codigogris},
  numbers=left, numbersep=6pt,
  backgroundcolor=\color{codigofondo},
  frame=single, rulecolor=\color{usfxazul!40},
  breaklines=true, showstringspaces=false,
  columns=fullflexible, keepspaces=true,
  xleftmargin=1.5em, framexleftmargin=1.5em,
  morekeywords={as,True,False,None},
}
\lstdefinestyle{salida}{
  basicstyle=\ttfamily\scriptsize,
  backgroundcolor=\color{black!85},
  basicstyle=\ttfamily\scriptsize\color{white},
  frame=none, numbers=none, columns=fullflexible, keepspaces=true,
}
\lstset{style=python}

% ---------------------------------------------------------------- Portada
\title[INT210 · Semana 1]{Fundamentos de IA, Machine Learning\\y la Jerarquía de Necesidades de Datos}
\subtitle{Machine Learning aplicado a la Ciberseguridad y al Diseño/Animación Digital}
\author[L. A. Molina Yampa]{Ing. Luis Antonio Molina Yampa}
\institute[USFX]{Universidad Mayor, Real y Pontificia de San Francisco Xavier de Chuquisaca\\
  INT210 -- Machine Learning}
\date{Semana 1}

\newcommand{\capa}[5]{%
  \fill[#4] (-#1,#2) -- (#1,#2) -- (#1-0.45,#2+0.9) -- (-#1+0.45,#2+0.9) -- cycle;
  \node[font=\scriptsize\bfseries, text=#5, align=center] at (0,#2+0.45) {#3};}
\newcommand{\ejemplo}[1]{\textcolor{usfxdorado}{\textbf{#1}}}

% =====================================================================
\begin{document}

\begin{frame}[plain]
  \titlepage
\end{frame}

% ---------------------------------------------------------------------
\begin{frame}{Objetivos de aprendizaje}
  Al finalizar la semana, el estudiante será capaz de:
  \medskip
  \begin{enumerate}
    \item \textbf{Diferenciar} IA, Machine Learning, Deep Learning y Estadística.
    \item \textbf{Representar} una muestra del mundo real como un vector de características $\mathbf{x}_i\in\mathbb{R}^d$.
    \item \textbf{Formalizar} el aprendizaje supervisado como la aproximación de una función $f:\mathcal{X}\to\mathcal{Y}$.
    \item \textbf{Clasificar} problemas según el tipo de aprendizaje (supervisado, no supervisado, semi-supervisado, refuerzo).
    \item \textbf{Explicar} por qué los datos (y no el modelo) son la base de todo proyecto de ML.
    \item \textbf{Interpretar} una matriz de confusión en un escenario de ciberseguridad.
  \end{enumerate}
\end{frame}

% ---------------------------------------------------------------------
\begin{frame}{Organización de la semana}
  \centering
  \begin{tabular}{>{\bfseries}l l l}
    \toprule
    Sesión & Modalidad & Producto \\
    \midrule
    Clase 1 (1.5 h) & Sesión teórica magistral & Conceptos y formalización \\
    Clase 2 (1.5 h) & Práctica guiada en Google Colab & Exploración del dataset Iris con Pandas \\
    Clase 3 (1.5 h) & Laboratorio evaluado & Clasificador heurístico + matriz de confusión \\
    Autónomo & Fuera de aula & Tarea + Quiz en Moodle \\
    \bottomrule
  \end{tabular}
  \vfill
  \begin{block}{Hilo conductor del curso}
    Cada concepto se ilustra con dos dominios de aplicación:
    \textbf{ciberseguridad} (tráfico de red, malware, intrusiones) y
    \textbf{diseño/animación digital} (píxeles, renders, fotogramas).
  \end{block}
\end{frame}

% =====================================================================
\section{IA, ML, DL y Estadística}
% =====================================================================

\begin{frame}{Delimitando los dominios}
  \begin{columns}[T]
    \begin{column}{0.47\textwidth}
      \centering
      \begin{tikzpicture}[scale=0.92]
        \draw[fill=usfxazul!12, draw=usfxazul, thick] (0,0) ellipse (3.1 and 2.25);
        \draw[fill=usfxazul!28, draw=usfxazul, thick] (0.45,-0.35) ellipse (2.25 and 1.6);
        \draw[fill=usfxazul!55, draw=usfxazul, thick] (0.9,-0.65) ellipse (1.25 and 0.9);
        \node[font=\bfseries\small, usfxazul] at (-0.2,1.75) {Inteligencia Artificial};
        \node[font=\bfseries\small, usfxazul] at (0.2,0.75) {Machine Learning};
        \node[font=\bfseries\small, white] at (0.9,-0.65) {Deep Learning};
        \node[font=\tiny, align=center, text=black!70] at (-2.05,-0.2) {Reglas\\Lógica difusa\\Heurísticas};
      \end{tikzpicture}
    \end{column}
    \begin{column}{0.51\textwidth}
      \small
      \begin{description}[IA]
        \item[IA] Superconjunto: cualquier técnica que permita a una computadora \emph{simular} razonamiento humano (incluye reglas fijas y heurísticas).
        \item[ML] Subcampo de la IA: algoritmos que \emph{aprenden patrones de los datos} sin ser programados explícitamente.
        \item[DL] Especialización del ML: redes neuronales multicapa que extraen \emph{jerarquías de características} de datos no estructurados (imagen, audio, video).
      \end{description}
    \end{column}
  \end{columns}
\end{frame}

\begin{frame}{Programación tradicional vs. Machine Learning}
  \centering
  \begin{tikzpicture}[
      caja/.style={draw=usfxazul, fill=usfxazul!8, rounded corners, minimum width=2.3cm, minimum height=0.8cm, font=\small},
      prog/.style={draw=usfxrojo, fill=usfxrojo!10, rounded corners, minimum width=2.6cm, minimum height=1.6cm, font=\small\bfseries, align=center},
      flecha/.style={-{Stealth[length=2.5mm]}, thick, usfxazul}]
    % Tradicional
    \node[font=\bfseries] at (-0.6,1.9) {Programación tradicional};
    \node[caja] (d1) at (-2.6,1.0) {Datos};
    \node[caja] (r1) at (-2.6,0.0) {Reglas};
    \node[prog] (p1) at (0.4,0.5) {Programa};
    \node[caja, fill=vp!15, draw=vp] (o1) at (3.2,0.5) {Respuestas};
    \draw[flecha] (d1) -- (p1); \draw[flecha] (r1) -- (p1); \draw[flecha] (p1) -- (o1);
    % ML
    \node[font=\bfseries] at (-0.6,-1.2) {Machine Learning};
    \node[caja] (d2) at (-2.6,-2.1) {Datos};
    \node[caja, fill=vp!15, draw=vp] (r2) at (-2.6,-3.1) {Respuestas};
    \node[prog] (p2) at (0.4,-2.6) {Algoritmo de\\aprendizaje};
    \node[caja, fill=usfxdorado!20, draw=usfxdorado] (o2) at (3.2,-2.6) {Reglas (modelo)};
    \draw[flecha] (d2) -- (p2); \draw[flecha] (r2) -- (p2); \draw[flecha] (p2) -- (o2);
  \end{tikzpicture}
  \vfill
  \begin{exampleblock}{Ejemplo -- Filtro de spam}
    \textbf{Tradicional:} \texttt{if "premio" in correo: spam}. \quad
    \textbf{ML:} se entrena con 10\,000 correos etiquetados y el modelo descubre por sí mismo qué palabras, remitentes y horarios delatan al spam.
  \end{exampleblock}
\end{frame}

\begin{frame}{Estadística vs. Machine Learning}
  \centering
  \renewcommand{\arraystretch}{1.3}
  \begin{tabular}{>{\bfseries}p{3.2cm} p{5.1cm} p{5.1cm}}
    \toprule
    & Estadística tradicional & Machine Learning \\
    \midrule
    Objetivo principal & Inferencia y explicación & Predicción y generalización \\
    Pregunta típica & ¿Es significativo este parámetro? & ¿Qué tan bien predice datos nuevos? \\
    Supuestos & Distribuciones previas estrictas & Pocos supuestos; validación empírica \\
    Escala & Muestras moderadas & Grandes volúmenes de datos \\
    Validación & $p$-valores, intervalos de confianza & Conjunto de prueba, validación cruzada \\
    \bottomrule
  \end{tabular}
  \vfill
  \begin{exampleblock}{Ejemplo -- mismo dato, dos preguntas}
    Con registros de un estudio de animación: el \textbf{estadístico} pregunta
    \emph{``¿el número de luces influye significativamente en el tiempo de render?''};
    el \textbf{ingeniero de ML} pregunta \emph{``¿cuántos segundos tardará el próximo fotograma?''}.
  \end{exampleblock}
\end{frame}

\begin{frame}{Analogía: el parque de diversiones}
  \begin{columns}[T]
    \begin{column}{0.32\textwidth}
      \begin{block}{IA: robot con manual}
        \small
        \emph{``Si ves una piedra roja, salta; si es azul, agáchate.''}\\[0.4em]
        Parece inteligente, pero ante una piedra \textbf{amarilla} se confunde: no está en sus reglas fijas.
      \end{block}
    \end{column}
    \begin{column}{0.32\textwidth}
      \begin{block}{ML: cachorro entrenable}
        \small
        Le mostramos 100 fotos de gatos y 100 de perros. Cada acierto recibe una ``galleta'' (premio matemático).\\[0.4em]
        Aprende \textbf{solo} a reconocer orejas y bigotes.
      \end{block}
    \end{column}
    \begin{column}{0.32\textwidth}
      \begin{block}{DL: detective con supergafas}
        \small
        Descompone la imagen en capas: \textbf{bordes} $\to$ \textbf{ojos} $\to$ \textbf{cara completa}.\\[0.4em]
        Procesa datos muy complejos como un cerebro de muchas capas.
      \end{block}
    \end{column}
  \end{columns}
  \vfill
  \begin{alertblock}{Pregunta para la clase}
    Un antivirus que bloquea archivos cuyo \emph{hash} está en una lista negra, ¿es IA, ML o DL? ¿Y uno que detecta malware \emph{nunca antes visto}?
  \end{alertblock}
\end{frame}

% =====================================================================
\section{Formalización matemática}
% =====================================================================

\begin{frame}{Representación vectorial de una muestra}
  Cada instancia del mundo real se transforma en un \textbf{vector de características}:
  \[
    \mathbf{x}_i = [x_{i,1},\, x_{i,2},\, \dots,\, x_{i,d}]^{T} \in \mathbb{R}^{d}
  \]
  \begin{itemize}
    \item $\mathbf{x}_i$: la $i$-ésima muestra \quad $\bullet$ \quad $d$: dimensionalidad (número de \emph{features}) \quad $\bullet$ \quad $x_{i,j}$: valor de la característica $j$.
  \end{itemize}
  \medskip
  \begin{columns}[T]
    \begin{column}{0.49\textwidth}
      \begin{exampleblock}{Ejemplo: conexión de red ($d=4$)}
        \small
        \[
          \mathbf{x}_{17} =
          \begin{bmatrix} 2.3 \\ 1540 \\ 12 \\ 443 \end{bmatrix}
          \begin{array}{l}
            \leftarrow \text{duración (s)} \\
            \leftarrow \text{bytes enviados} \\
            \leftarrow \text{nº de paquetes} \\
            \leftarrow \text{puerto destino}
          \end{array}
        \]
      \end{exampleblock}
    \end{column}
    \begin{column}{0.49\textwidth}
      \begin{exampleblock}{Ejemplo: flor Iris ($d=4$)}
        \small
        \[
          \mathbf{x}_{1} =
          \begin{bmatrix} 5.1 \\ 3.5 \\ 1.4 \\ 0.2 \end{bmatrix}
          \begin{array}{l}
            \leftarrow \text{largo del sépalo (cm)} \\
            \leftarrow \text{ancho del sépalo (cm)} \\
            \leftarrow \text{largo del pétalo (cm)} \\
            \leftarrow \text{ancho del pétalo (cm)}
          \end{array}
        \]
      \end{exampleblock}
    \end{column}
  \end{columns}
\end{frame}

\begin{frame}{El dataset como matriz}
  Con $n$ muestras apiladas obtenemos la \textbf{matriz de diseño} $\mathbf{X}\in\mathbb{R}^{n\times d}$ y el vector de etiquetas $\mathbf{y}$:
  \[
    \mathbf{X} =
    \begin{bmatrix}
      x_{1,1} & x_{1,2} & \cdots & x_{1,d} \\
      x_{2,1} & x_{2,2} & \cdots & x_{2,d} \\
      \vdots  & \vdots  & \ddots & \vdots  \\
      x_{n,1} & x_{n,2} & \cdots & x_{n,d}
    \end{bmatrix},
    \qquad
    \mathbf{y} = \begin{bmatrix} y_1 \\ y_2 \\ \vdots \\ y_n \end{bmatrix}
  \]
  \begin{exampleblock}{Ejemplo: Iris}
    \centering\small
    \begin{tabular}{c cccc c}
      \toprule
      $i$ & sépalo largo & sépalo ancho & pétalo largo & pétalo ancho & $y_i$ \\
      \midrule
      1   & 5.1 & 3.5 & 1.4 & 0.2 & 0 (setosa) \\
      51  & 7.0 & 3.2 & 4.7 & 1.4 & 1 (versicolor) \\
      101 & 6.3 & 3.3 & 6.0 & 2.5 & 2 (virginica) \\
      \bottomrule
    \end{tabular}\\[0.3em]
    Iris completo: $n=150$, $d=4$ $\;\Rightarrow\;$ $\mathbf{X}\in\mathbb{R}^{150\times 4}$.
  \end{exampleblock}
\end{frame}

\begin{frame}{La función de mapeo predictivo}
  El aprendizaje supervisado busca aproximar una función
  \[
    f:\mathcal{X}\to\mathcal{Y}, \qquad \hat{y}_i = f(\mathbf{x}_i)
  \]
  que minimice el error empírico sobre los datos de entrenamiento.
  \medskip
  \begin{columns}[T]
    \begin{column}{0.49\textwidth}
      \begin{block}{Clasificación}
        $\mathcal{Y}$ es un conjunto \textbf{discreto}:
        \[ \mathcal{Y}=\{0,1\} \quad\text{o}\quad \mathcal{Y}=\{-1,+1\} \]
        \small\ejemplo{Ejemplo:} ¿el tráfico es \emph{seguro} (0) o un \emph{ataque DDoS} (1)?
      \end{block}
    \end{column}
    \begin{column}{0.49\textwidth}
      \begin{block}{Regresión}
        $\mathcal{Y}$ es \textbf{continuo}:
        \[ \mathcal{Y}=\mathbb{R} \]
        \small\ejemplo{Ejemplo:} ¿cuántos segundos tardará en renderizarse un fotograma? \; $\hat{y}=42.7$ s
      \end{block}
    \end{column}
  \end{columns}
  \medskip
  \begin{alertblock}{¿Cómo se mide el error? -- Función de pérdida}
    \small
    Clasificación (0-1): $\mathcal{L}(y,\hat{y}) = \mathbf{1}[y\neq\hat{y}]$
    \qquad
    Regresión (cuadrática): $\mathcal{L}(y,\hat{y}) = (y-\hat{y})^2$
  \end{alertblock}
\end{frame}

\begin{frame}{Ejemplo resuelto: calculando la pérdida}
  Un modelo predice el tiempo de render (s) de 4 fotogramas:
  \medskip
  \begin{columns}[c]
    \begin{column}{0.48\textwidth}
      \centering
      \begin{tabular}{cccc}
        \toprule
        Fotograma & $y_i$ & $\hat{y}_i$ & $(y_i-\hat{y}_i)^2$ \\
        \midrule
        1 & 40 & 38 & 4 \\
        2 & 55 & 60 & 25 \\
        3 & 32 & 31 & 1 \\
        4 & 70 & 64 & 36 \\
        \midrule
        \multicolumn{3}{r}{\textbf{Suma}} & \textbf{66} \\
        \bottomrule
      \end{tabular}
    \end{column}
    \begin{column}{0.48\textwidth}
      Error cuadrático medio (MSE):
      \[
        \text{MSE} = \frac{1}{n}\sum_{i=1}^{n}(y_i-\hat{y}_i)^2 = \frac{66}{4} = 16.5\ \text{s}^2
      \]
      Raíz del MSE:
      \[
        \text{RMSE} = \sqrt{16.5} \approx 4.06\ \text{s}
      \]
      \small\emph{Interpretación:} en promedio el modelo se equivoca en unos 4 segundos por fotograma.
    \end{column}
  \end{columns}
\end{frame}

\begin{frame}{Tipos de aprendizaje}
  \small
  \renewcommand{\arraystretch}{1.35}
  \begin{tabular}{>{\bfseries}p{2.6cm} p{4.2cm} p{6.6cm}}
    \toprule
    Tipo & Datos disponibles & Ejemplo en nuestros dominios \\
    \midrule
    Supervisado & $\mathcal{D}=\{(\mathbf{x}_i,y_i)\}_{i=1}^n$ \newline minimiza $\mathcal{L}(y_i,f(\mathbf{x}_i))$ & Clasificar correos como \emph{phishing} / legítimo a partir de miles de ejemplos etiquetados. \\
    No supervisado & $\mathcal{D}=\{\mathbf{x}_i\}_{i=1}^n$ \newline (sin etiquetas) & Agrupar (clustering) sesiones de usuarios por comportamiento; detectar un grupo anómalo. Agrupar paletas de color de un portafolio. \\
    Semi-supervisado & Pocos datos etiquetados + muchos sin etiquetar & 200 binarios etiquetados por un analista + 50\,000 sin revisar para mejorar un detector de malware a bajo costo. \\
    Por refuerzo & Agente + entorno + recompensas & Un personaje animado aprende a caminar; un agente aprende a responder ataques en un simulador de red. \\
    \bottomrule
  \end{tabular}
\end{frame}

\begin{frame}{Actividad rápida: ¿qué tipo de problema es?}
  Para cada caso indique: \textbf{(a)} tipo de aprendizaje y \textbf{(b)} clasificación o regresión (si aplica).
  \medskip
  \begin{enumerate}
    \item Predecir cuántos megabytes ocupará un video exportado a partir de su duración y resolución.
    \item Separar 1 millón de registros de \emph{logs} en grupos sin conocer de antemano qué representan.
    \item Decidir si un inicio de sesión es fraudulento usando un historial de accesos marcados por el SOC.
    \item Entrenar un bot que aprende a jugar un videojuego maximizando su puntaje.
  \end{enumerate}
  \vfill
  \pause
  \begin{exampleblock}{Respuestas}
    \small
    1.~Supervisado -- regresión \quad
    2.~No supervisado -- clustering \quad
    3.~Supervisado -- clasificación \quad
    4.~Refuerzo
  \end{exampleblock}
\end{frame}

% =====================================================================
\section{Jerarquía de necesidades de datos}
% =====================================================================

\begin{frame}{La pirámide de Monica Rogati}
  \begin{columns}[c]
    \begin{column}{0.55\textwidth}
      \centering
      \begin{tikzpicture}[scale=0.78]
        \capa{4.2}{0.0}{Recolectar\\{\tiny sensores, logs de red, APIs}}{usfxazul!25}{black}
        \capa{3.75}{0.9}{Mover / Almacenar\\{\tiny ingesta, pipelines ETL}}{usfxazul!40}{black}
        \capa{3.3}{1.8}{Explorar / Transformar\\{\tiny limpieza, anomalías}}{usfxazul!60}{white}
        \capa{2.85}{2.7}{Agregar / Etiquetar\\{\tiny métricas, features}}{usfxazul!75}{white}
        \capa{2.4}{3.6}{Aprender / Optimizar\\{\tiny A/B testing, reglas}}{usfxazul!90}{white}
        \fill[usfxdorado] (-1.95,4.5) -- (1.95,4.5) -- (0,5.7) -- cycle;
        \node[font=\scriptsize\bfseries] at (0,4.9) {IA / ML};
        \draw[thick, usfxrojo, |-|] (4.4,0) -- (4.4,2.7) node[midway, right, font=\scriptsize, align=left, text=usfxrojo] {70\%--80\%\\del tiempo};
      \end{tikzpicture}
    \end{column}
    \begin{column}{0.43\textwidth}
      \small
      \begin{itemize}
        \item No se puede modelar sin una \textbf{base robusta} de datos.
        \item Reis y Housley (2022): entre el \textbf{70\% y 80\%} del tiempo de un científico de datos se consume en las capas inferiores.
        \item El modelo es la \emph{punta} visible; la ingeniería de datos es el iceberg.
      \end{itemize}
      \begin{alertblock}{Antipatrón}
        \footnotesize Entrenar un clasificador de malware o un generador de arte sin pipelines estables ni limpieza previa: costoso e inviable.
      \end{alertblock}
    \end{column}
  \end{columns}
\end{frame}

\begin{frame}{Ejemplo: la pirámide en un SOC de un estudio de animación}
  \small
  \renewcommand{\arraystretch}{1.3}
  \begin{tabular}{>{\bfseries}p{3.2cm} p{10.2cm}}
    \toprule
    Capa & Qué se hace en la práctica \\
    \midrule
    Recolectar & Capturar logs del firewall, NetFlow de la granja de render, eventos de inicio de sesión. \\
    Mover / Almacenar & Enviar los logs a un repositorio central (p. ej. un SIEM) con un pipeline ETL confiable. \\
    Explorar / Transformar & Eliminar duplicados, normalizar zonas horarias, corregir IPs mal formadas. \\
    Agregar / Etiquetar & Calcular ``bytes por minuto por host''; el analista marca incidentes reales como \emph{ataque}. \\
    Aprender / Optimizar & Probar primero reglas simples (umbral de bytes) y medir su efecto. \\
    \textcolor{usfxdorado}{IA / ML} & Recién ahora: entrenar un modelo que detecte exfiltración de renders confidenciales. \\
    \bottomrule
  \end{tabular}
\end{frame}

% =====================================================================
\section{Práctica guiada: explorando Iris con Pandas}
% =====================================================================

\begin{frame}{El dataset Iris}
  \begin{columns}[c]
    \begin{column}{0.45\textwidth}
      \begin{itemize}
        \item $n = 150$ flores, 50 por especie.
        \item $d = 4$ características numéricas (cm).
        \item 3 clases: \emph{setosa}, \emph{versicolor}, \emph{virginica}.
        \item Incluido en \texttt{scikit-learn}: no requiere descarga manual.
      \end{itemize}
      \medskip
      \begin{block}{¿Por qué Iris?}
        \small Es pequeño, limpio y visualizable: ideal para aprender el \emph{flujo de trabajo} antes de usar datos reales de red.
      \end{block}
    \end{column}
    \begin{column}{0.53\textwidth}
      \centering\footnotesize
      \begin{tabular}{l ccc}
        \toprule
        & \multicolumn{3}{c}{Largo del pétalo (cm)} \\
        \cmidrule(lr){2-4}
        Especie & mín & máx & media \\
        \midrule
        setosa     & 1.0 & 1.9 & 1.46 \\
        versicolor & 3.0 & 5.1 & 4.26 \\
        virginica  & 4.5 & 6.9 & 5.55 \\
        \midrule
        & \multicolumn{3}{c}{Ancho del pétalo (cm)} \\
        \cmidrule(lr){2-4}
        setosa     & 0.1 & 0.6 & 0.25 \\
        versicolor & 1.0 & 1.8 & 1.33 \\
        virginica  & 1.4 & 2.5 & 2.03 \\
        \bottomrule
      \end{tabular}\\[0.4em]
      \emph{Observación:} ¡las setosas no se solapan con las demás!
    \end{column}
  \end{columns}
\end{frame}

\begin{frame}[fragile]{Paso 1: Carga de datos}
  \begin{lstlisting}
import pandas as pd
from sklearn.datasets import load_iris

iris_raw = load_iris()                                  # Carga del dataset Iris
df = pd.DataFrame(data=iris_raw.data, columns=iris_raw.feature_names)
df['species'] = iris_raw.target                         # Etiquetas 0, 1, 2

species_map = {0: 'setosa', 1: 'versicolor', 2: 'virginica'}
df['species_name'] = df['species'].map(species_map)     # Mapeo semántico
df.head()
  \end{lstlisting}
  \textbf{Salida:}
  \begin{lstlisting}[style=salida]
   sepal length (cm)  sepal width (cm)  petal length (cm)  petal width (cm)  species species_name
0                5.1               3.5                1.4               0.2        0       setosa
1                4.9               3.0                1.4               0.2        0       setosa
2                4.7               3.2                1.3               0.2        0       setosa
3                4.6               3.1                1.5               0.2        0       setosa
4                5.0               3.6                1.4               0.2        0       setosa
  \end{lstlisting}
\end{frame}

\begin{frame}{Paso 1: ¿qué hace cada línea?}
  \small
  \begin{description}[\texttt{df['species\_name']}]
    \item[\texttt{import pandas as pd}] Librería estándar para estructuras tabulares (\emph{DataFrames}).
    \item[\texttt{load\_iris()}] Devuelve de forma determinista las 150 muestras con sus 4 atributos.
    \item[\texttt{pd.DataFrame(...)}] Crea la tabla asignando nombres a cada columna (\emph{feature}).
    \item[\texttt{df['species']}] Agrega la etiqueta real $y_i \in \{0,1,2\}$.
    \item[\texttt{df['species\_name']}] Traduce números a nombres legibles con \texttt{.map()}.
  \end{description}
  \vfill
  \begin{block}{Conexión con la teoría}
    Las 4 primeras columnas forman $\mathbf{X}\in\mathbb{R}^{150\times4}$; la columna \texttt{species} es el vector $\mathbf{y}$.
  \end{block}
\end{frame}

\begin{frame}[fragile]{Paso 2: Consultas con \texttt{query()}}
  \textbf{Objetivo:} encontrar setosas con pétalos muy delgados (ancho $\le 0.2$ cm), como un \emph{filtro de umbral bajo}.
  \begin{lstlisting}
filtro_setosa = df.query("species_name == 'setosa' and `petal width (cm)` <= 0.2")
print(f"Total de muestras encontradas: {len(filtro_setosa)}")
  \end{lstlisting}
  \begin{lstlisting}[style=salida]
Total de muestras encontradas: 34
  \end{lstlisting}
  \medskip
  \begin{itemize}
    \item Los nombres de columna con espacios se encierran entre \textbf{comillas invertidas} (\texttt{\`{}...\`{}}).
    \item \texttt{and} exige que \textbf{ambas} condiciones se cumplan simultáneamente.
  \end{itemize}
  \begin{exampleblock}{Analogía en ciberseguridad}
    \small\texttt{logs.query("puerto == 22 and `intentos\_fallidos` >= 5")} \; $\to$ posibles ataques de fuerza bruta por SSH.
  \end{exampleblock}
\end{frame}

\begin{frame}[fragile]{Paso 3: Correlación de Pearson}
  \begin{columns}[c]
    \begin{column}{0.5\textwidth}
\begin{lstlisting}
cols = ['petal length (cm)',
        'petal width (cm)']
corr = df[cols].corr(method='pearson')
print(corr)
\end{lstlisting}
\begin{lstlisting}[style=salida]
                   petal length  petal width
petal length (cm)      1.000000     0.962865
petal width (cm)       0.962865     1.000000
\end{lstlisting}
      \small
      \[
        r = \frac{\sum_i (x_i-\bar{x})(y_i-\bar{y})}{\sqrt{\sum_i (x_i-\bar{x})^2}\sqrt{\sum_i (y_i-\bar{y})^2}} \in [-1,1]
      \]
      $r\approx0.96$: relación lineal positiva \textbf{muy fuerte}.
    \end{column}
    \begin{column}{0.48\textwidth}
      \centering
      \begin{tikzpicture}
        \begin{axis}[
            width=6.6cm, height=5.6cm,
            xlabel={Largo del pétalo (cm)}, ylabel={Ancho del pétalo (cm)},
            label style={font=\scriptsize}, tick label style={font=\tiny},
            legend style={font=\tiny, at={(0.97,0.03)}, anchor=south east},
            scatter/classes={
              0={mark=*, usfxazul, mark size=1.3pt},
              1={mark=square*, usfxdorado, mark size=1.3pt},
              2={mark=triangle*, usfxrojo, mark size=1.6pt}},
            xmin=0.5, xmax=7.3, ymin=0, ymax=2.7]
          \addplot[scatter, only marks, scatter src=explicit symbolic]
            table[x=largo, y=ancho, meta=clase] {iris_petalos.dat};
          \legend{setosa, versicolor, virginica}
          \draw[dashed, thick, black!70] (axis cs:2.0,0) -- (axis cs:2.0,2.7);
          \node[font=\tiny, anchor=west] at (axis cs:2.05,2.35) {umbral $=2.0$};
        \end{axis}
      \end{tikzpicture}
    \end{column}
  \end{columns}
\end{frame}

% =====================================================================
\section{Laboratorio evaluado: heurística vs. modelo}
% =====================================================================

\begin{frame}{El desafío del laboratorio}
  \begin{enumerate}
    \item \textbf{Cargar} el dataset Iris (estructura de la Clase 2).
    \item \textbf{Construir un modelo heurístico} con \texttt{if-elif-else} que prediga \emph{setosa} vs.\ \emph{no-setosa} usando un umbral sobre el pétalo.
    \item \textbf{Comparar} con un modelo algorítmico: $k$-vecinos más cercanos ($k$-NN) con $k=3$.
    \item \textbf{Generar la matriz de confusión} y defender sus implicaciones en un escenario de ciberseguridad.
  \end{enumerate}
  \vfill
  \begin{alertblock}{Escenario simulado}
    \centering
    \textbf{Setosa} $=$ Alerta de intrusión crítica (\emph{malware}) \qquad
    \textbf{No-setosa} $=$ Tráfico de red legítimo
  \end{alertblock}
\end{frame}

\begin{frame}[fragile]{Modelo heurístico (hecho a mano)}
  \begin{lstlisting}
import numpy as np
from sklearn.metrics import confusion_matrix

def clasificador_manual(row):
    # Regla basada en el análisis visual: las setosas tienen pétalos muy cortos
    if row['petal length (cm)'] < 2.0:
        return 0  # Setosa  (=> Malware)
    else:
        return 1  # No-Setosa (=> Tráfico legítimo)

df['prediccion_heuristica'] = df.apply(clasificador_manual, axis=1)
df['label_binario'] = df['species'].apply(lambda x: 0 if x == 0 else 1)

conf_matrix = confusion_matrix(df['label_binario'], df['prediccion_heuristica'])
print("Matriz de Confusión Obtenida:")
print(conf_matrix)
  \end{lstlisting}
  \begin{lstlisting}[style=salida]
Matriz de Confusión Obtenida:
[[ 50   0]
 [  0 100]]
  \end{lstlisting}
\end{frame}

\begin{frame}{Leyendo la matriz de confusión}
  \begin{columns}[c]
    \begin{column}{0.5\textwidth}
      \centering
      \begin{tikzpicture}[cel/.style={minimum width=2.2cm, minimum height=1.25cm, draw=white, very thick, align=center, font=\small}]
        \node[cel, fill=vp!80, text=white]  (a) at (0,0)     {\textbf{VP}\\50};
        \node[cel, fill=fn!80, text=white]  (b) at (2.2,0)   {\textbf{FN}\\0};
        \node[cel, fill=fp!80, text=white]  (c) at (0,-1.25) {\textbf{FP}\\0};
        \node[cel, fill=vp!45]              (d) at (2.2,-1.25){\textbf{VN}\\100};
        \node[font=\scriptsize, above=0.05cm of a] {Pred.\ Malware};
        \node[font=\scriptsize, above=0.05cm of b] {Pred.\ Legítimo};
        \node[font=\scriptsize, rotate=90, align=center] at ($(a.west)+(-0.4,0)$) {Real\\Malware};
        \node[font=\scriptsize, rotate=90, align=center] at ($(c.west)+(-0.4,0)$) {Real\\Legítimo};
      \end{tikzpicture}
    \end{column}
    \begin{column}{0.48\textwidth}
      \small
      \begin{itemize}
        \item Filas $=$ clase \textbf{real}; columnas $=$ clase \textbf{predicha}.
        \item \texttt{scikit-learn} ordena las etiquetas: la clase \texttt{0} (setosa $=$ malware, nuestra clase \emph{positiva}) aparece primero.
        \item Resultado: \textbf{clasificación perfecta}. El umbral de 2.0 cm separa totalmente a las setosas.
      \end{itemize}
      \begin{alertblock}{¡Cuidado!}
        \footnotesize En la mayoría de los textos la clase positiva es \texttt{1} y aparece abajo a la derecha. Verifique siempre el orden de las etiquetas.
      \end{alertblock}
    \end{column}
  \end{columns}
\end{frame}

\begin{frame}{Defensa analítica: los cuatro extremos de la matriz}
  \small
  \begin{columns}[T]
    \begin{column}{0.49\textwidth}
      \begin{exampleblock}{Verdadero Positivo (VP)}
        Predice malware y \textbf{es} malware. El firewall intercepta la amenaza y protege los activos del estudio de diseño.
      \end{exampleblock}
      \begin{block}{Falso Positivo (FP) -- Error Tipo I}
        Predice malware pero era \textbf{tráfico legítimo} de renders 3D. Se corta la conexión con la granja de servidores: retrasos y falsas alarmas que saturan al SOC.
      \end{block}
    \end{column}
    \begin{column}{0.49\textwidth}
      \begin{exampleblock}{Verdadero Negativo (VN)}
        Predice tráfico legítimo y \textbf{lo es}. Eficiencia operativa: los animadores renderizan sin bloqueos.
      \end{exampleblock}
      \begin{alertblock}{Falso Negativo (FN) -- Error Tipo II}
        Predice legítimo pero era \textbf{malware}. Un ransomware entra en silencio, cifra los servidores y compromete la propiedad intelectual (modelos 3D bajo NDA). \textbf{El error más peligroso.}
      \end{alertblock}
    \end{column}
  \end{columns}
\end{frame}

\begin{frame}{Ejemplo con errores: un IDS en producción}
  Un sistema de detección analiza \textbf{1\,000 paquetes}; 50 son maliciosos.
  \medskip
  \begin{columns}[c]
    \begin{column}{0.42\textwidth}
      \centering
      \begin{tikzpicture}[cel/.style={minimum width=2.0cm, minimum height=1.15cm, draw=white, very thick, align=center, font=\small}]
        \node[cel, fill=vp!80, text=white]   (a) at (0,0)     {\textbf{VP}\\45};
        \node[cel, fill=fn!80, text=white]   (b) at (2.0,0)   {\textbf{FN}\\5};
        \node[cel, fill=fp!80, text=white]   (c) at (0,-1.15) {\textbf{FP}\\30};
        \node[cel, fill=vp!45]               (d) at (2.0,-1.15){\textbf{VN}\\920};
        \node[font=\scriptsize, above=0.05cm of a] {Pred.\ Ataque};
        \node[font=\scriptsize, above=0.05cm of b] {Pred.\ Normal};
        \node[font=\scriptsize, rotate=90, align=center] at ($(a.west)+(-0.4,0)$) {Real\\Ataque};
        \node[font=\scriptsize, rotate=90, align=center] at ($(c.west)+(-0.4,0)$) {Real\\Normal};
      \end{tikzpicture}
    \end{column}
    \begin{column}{0.56\textwidth}
      \small
      \begin{align*}
        \text{Exactitud} &= \tfrac{VP+VN}{\text{Total}} = \tfrac{45+920}{1000} = 96.5\% \\
        \text{Precisión} &= \tfrac{VP}{VP+FP} = \tfrac{45}{75} = 60.0\% \\
        \text{Sensibilidad (Recall)} &= \tfrac{VP}{VP+FN} = \tfrac{45}{50} = 90.0\%
      \end{align*}
    \end{column}
  \end{columns}
  \medskip
  \begin{alertblock}{Lección}
    \small Una exactitud de 96.5\% suena excelente, pero \textbf{5 ataques pasaron sin ser detectados} y 4 de cada 10 alertas son falsas. En datos desbalanceados la exactitud \emph{engaña}.
  \end{alertblock}
\end{frame}

\begin{frame}[fragile]{Modelo algorítmico: $k$-NN con $k=3$}
  \begin{columns}[T]
    \begin{column}{0.55\textwidth}
\begin{lstlisting}
from sklearn.neighbors import KNeighborsClassifier
from sklearn.model_selection import train_test_split

X = df[iris_raw.feature_names]
y = df['species']
X_tr, X_te, y_tr, y_te = train_test_split(
    X, y, test_size=0.3,
    random_state=42, stratify=y)

knn = KNeighborsClassifier(n_neighbors=3)
knn.fit(X_tr, y_tr)
print(knn.score(X_te, y_te))
print(confusion_matrix(y_te, knn.predict(X_te)))
\end{lstlisting}
\begin{lstlisting}[style=salida]
0.9555555555555556
[[15  0  0]
 [ 0 15  0]
 [ 0  2 13]]
\end{lstlisting}
    \end{column}
    \begin{column}{0.43\textwidth}
      \small
      \begin{block}{Idea de $k$-NN}
        Para clasificar una flor nueva, se buscan las $k=3$ flores más \textbf{cercanas} (distancia euclidiana) y se vota:
        \[ d(\mathbf{x},\mathbf{z})=\sqrt{\textstyle\sum_{j=1}^{d}(x_j-z_j)^2} \]
      \end{block}
      \begin{itemize}
        \item Exactitud en prueba: \textbf{95.6\%} (3 clases).
        \item 2 virginicas confundidas con versicolor: las clases que \emph{sí} se solapan.
        \item La heurística sólo resuelve el problema \emph{binario} fácil.
      \end{itemize}
    \end{column}
  \end{columns}
\end{frame}

% =====================================================================
\section{Trabajo autónomo y cierre}
% =====================================================================

\begin{frame}{Tarea para la casa}
  \begin{block}{Instrucciones}
    \begin{enumerate}
      \item Modifique \texttt{clasificador\_manual} para usar \texttt{petal width (cm)} en lugar del largo.
      \item Encuentre \textbf{iterativamente} el umbral que lleve los \textbf{Falsos Negativos a cero}, aunque aumenten levemente los Falsos Positivos.
      \item Argumente matemáticamente por qué esta parametrización es preferible en ciberseguridad reactiva.
    \end{enumerate}
  \end{block}
  \begin{exampleblock}{Pista}
    \small
    Revise la tabla de estadísticos: el ancho máximo de una setosa es 0.6 cm y el mínimo de una versicolor es 1.0 cm.
    ¿Qué pasa con la regla \texttt{ancho < 0.6}? ¿Y con \texttt{ancho <= 0.6}?
    Relacione su respuesta con la \textbf{sensibilidad} $\frac{VP}{VP+FN}$.
  \end{exampleblock}
\end{frame}

\begin{frame}{Quiz de Moodle: preguntas de muestra}
  \small
  \begin{enumerate}
    \item ¿Qué subcampo de la IA entrena algoritmos para aprender patrones de los datos sin ser programados explícitamente?
      \begin{enumerate}[A.]
        \item Inteligencia Artificial General \quad B. Machine Learning \quad C. Sistemas expertos \quad D. PLN
      \end{enumerate}
    \item En detección de intrusiones, un \textbf{Falso Negativo} (Error Tipo II) es\ldots
      \begin{enumerate}[A.]
        \item Clasificar tráfico legítimo como intrusión.
        \item Clasificar tráfico malicioso como acceso seguro y legítimo.
      \end{enumerate}
    \item En $\mathbf{x}=[x_1,x_2,\dots,x_d]$, ¿qué representa $d$?
      \begin{enumerate}[A.]
        \item Número de muestras \quad B. La etiqueta \quad C. La dimensionalidad \quad D. Coeficiente de Pearson
      \end{enumerate}
  \end{enumerate}
  \vfill
  \centering\footnotesize\textcolor{codigogris}{El quiz completo (formato Aiken) se publica en Moodle.}
\end{frame}

\begin{frame}{Resumen de la semana}
  \begin{columns}[T]
    \begin{column}{0.49\textwidth}
      \begin{block}{Conceptos clave}
        \small
        \begin{itemize}
          \item IA $\supset$ ML $\supset$ DL.
          \item ML prioriza \textbf{predicción}; Estadística, \textbf{inferencia}.
          \item Muestra $=$ vector $\mathbf{x}_i\in\mathbb{R}^d$; dataset $=$ matriz $\mathbf{X}\in\mathbb{R}^{n\times d}$.
          \item Supervisado: aprender $f:\mathcal{X}\to\mathcal{Y}$ minimizando $\mathcal{L}$.
        \end{itemize}
      \end{block}
    \end{column}
    \begin{column}{0.49\textwidth}
      \begin{block}{Ideas para llevar}
        \small
        \begin{itemize}
          \item Los datos son la base de la pirámide: sin ellos no hay ML.
          \item Explorar antes de modelar (\texttt{query}, \texttt{corr}, gráficos).
          \item En seguridad, el \textbf{FN} es el error más caro.
          \item La exactitud sola puede engañar.
        \end{itemize}
      \end{block}
    \end{column}
  \end{columns}
  \vfill
  \begin{alertblock}{Próxima semana}
    \centering Continuamos con la Semana 2 del curso. ¡Traer el notebook de esta semana funcionando!
  \end{alertblock}
\end{frame}

\begin{frame}{Referencias bibliográficas}
  \small
  \begin{thebibliography}{9}
    \setbeamertemplate{bibliography item}[book]
    \bibitem{grus} Grus, J. (2019).
      \newblock \emph{Data Science from Scratch: First Principles with Python} (2nd ed.).
      \newblock O'Reilly Media.
    \bibitem{muller} Müller, A. C., \& Guido, S. (2017).
      \newblock \emph{Introduction to Machine Learning with Python: A Guide for Data Scientists}.
      \newblock O'Reilly Media.
    \bibitem{reis} Reis, J., \& Housley, M. (2022).
      \newblock \emph{Fundamentals of Data Engineering: Plan and Build Robust Data Systems}.
      \newblock O'Reilly Media.
    \bibitem{jamison} Jamison, P. (2023).
      \newblock \emph{The (Absolute) Beginner's Guide to AI}.
      \newblock Paul Jamison Publishing.
  \end{thebibliography}
\end{frame}

\begin{frame}[plain]
  \centering
  \vfill
  {\Huge\bfseries\color{usfxazul} ¡Gracias!}\\[1em]
  {\large ¿Preguntas?}\\[2em]
  \small INT210 -- Machine Learning \quad·\quad USFX
  \vfill
\end{frame}

\end{document}
